% Addition to chapters/09-zero-geometry.tex.
% Label prefix reserved for this continuation: lerchcont:
% Requires only the theorem environments and amsmath already in the manuscript.
% Add the bibliography item in bibliography-item.tex before compiling.
\section{Lerch endpoint singularities and zero bifurcations}
\label{lerchcont:section}
The continuation \cite{lerchcont2026} gives an exact description of the
Stieltjes endpoint, a low-order obstruction to all-branch monotonicity, and
the first sharp deficient classical zero count. Its full analytic proofs and
integer interval certificates accompany the standalone article. No numerical
root list is used as a substitute for a completeness argument.

Use the notation already established in this chapter:
\[
 F_{n,k}^{\rho}(a)=(-1)^{n+k}\mathcal D_{n,k}^{\rho}(a)
 =n![t^n](1+t)_k\Phi(\rho,k+1+t,a),\qquad (1+t)_k=k!P_k(t).
\]
At $\rho=1$, this is $(-1)^{n+k}\gamma_n^{(k)}(a)$.

\subsection{The universal logarithmic block}
Define new gamma-moment polynomials $R_m$, distinct from $Q_n$, by
\[
 \Gamma(1-t)e^{Xt}=\sum_{m\ge0}R_m(X)\frac{t^m}{m!}.
\]
For $r\ne k$, put
\[
 A_{n,k,r}(a)=n![t^n](1+t)_k\zeta(k+1-r+t,a),\qquad
 C_{n,k}(a)=n![t^n]\left(P_k(t)\zeta(1+t,a)-\frac1t\right).
\]
\begin{theorem}[Exact endpoint decomposition]\label{lerchcont:endpoint}
For $n\ge0$, $k\ge1$, $a>0$, and $0<\mu<2\pi$,
\begin{equation}\label{lerchcont:decomposition}
\begin{aligned}
 e^{-a\mu}F_{n,k}^{e^{-\mu}}(a)
 ={}&\sum_{\substack{r\ge0\\r\ne k}}\frac{(-\mu)^r}{r!}A_{n,k,r}(a)\\
 &+(-1)^k\mu^k\left(C_{n,k}(a)
                 -\frac{R_{n+1}(\log\mu)}{n+1}\right).
\end{aligned}
\end{equation}
The sum is an analytic power series at $\mu=0$. All nonanalytic terms occur
in the single displayed $\mu^k$ times logarithmic polynomial.
\end{theorem}
\begin{proof}
Apply the classical Erd\'elyi expansion for
$e^{-a\mu}\Phi(e^{-\mu},s,a)$, substitute $s=k+1+t$, and multiply by
$(1+t)_k$. The identity
\[
 (1+t)_k\Gamma(-k-t)=(-1)^{k+1}\Gamma(1-t)/t
\]
makes its pole cancel the $r=k$ Hurwitz term. Extracting $n![t^n]$ proves
\eqref{lerchcont:decomposition}. The classical convergent expansion justifies
coefficient extraction and differentiation on compact positive $a$ intervals.
\end{proof}

In particular, at $a=1$ this gives finite-part identities for order derivatives
of $\operatorname{Li}_{k+1+t}(e^{-\mu})$. With $Y=X+\gamma$,
\[
 R_3(X)=Y^3+3\zeta(2)Y+2\zeta(3),\qquad
 C_{2,1}(a)=\gamma_2(a)-2\gamma_1(a).
\]
The full coefficient recursion and concrete weight-two and weight-three
order-derivative limits are given in \cite{lerchcont2026}.

\begin{corollary}[Sharp endpoint smoothness]\label{lerchcont:regularity}
For fixed $n,k,a$, the function $\rho\mapsto F_{n,k}^{\rho}(a)$ is
$C^{k-1}$ but not $C^k$ at one. Writing $L=\log(1/\mu)$,
\[
 \partial_\rho^k F_{n,k}^{\rho}(a)
 \sim\frac{(-1)^n k!}{n+1}L^{n+1}.
\]
If $\alpha$ is a simple zero of $F_{n,1}^{1}$, its nearby branch satisfies
\[
 a'(\rho)\sim
 \frac{(-1)^{n+1}L^{n+1}}
 {(n+1)\partial_aF_{n,1}^{1}(\alpha)}.
\]
\end{corollary}
The first conclusion follows by differentiating the exact convergent
decomposition, and the second by implicit differentiation at a simple zero.
Thus uniform continuity of the family must not be confused with endpoint
smoothness of the same order.

\subsection{A monotonicity dichotomy}
\begin{theorem}[A unique minimum at index two]\label{lerchcont:n2}
For every $\rho\in[0,1]$, $F_{2,1}^{\rho}$ has exactly two positive zeros,
both simple. The upper zero is strictly decreasing in $\rho$. The lower zero
$a_-(\rho)$ has a unique nondegenerate interior minimum, is strictly decreasing
before it, and strictly increasing afterwards. Moreover,
\[
 a_-(0)=1,\qquad
 a_-'(0)=-\frac{\log2(2-\log2)}8<0,\qquad a_-(1)>1,
 \qquad a_-'(\rho)\longrightarrow+\infty\quad(\rho\uparrow1).
\]
\end{theorem}
\begin{proof}[Proof mechanism]
The exact certificate gives $F_{2,1}^{1}(13/10)<0$ and $F_{2,1}^{1}(1)>0$.
In the convergent series
\[
 F_{2,1}^{\rho}(a)=\sum_{m\ge0}\rho^m
 \frac{\log(a+m)(\log(a+m)-2)}{(a+m)^2},
\]
the coefficient signs change once for $1<a<e^2$, and at most twice for
$0<a<1$. The one-change argument propagates the negative separator to every
$\rho$. Endpoint signs and the global bound give exactly two simple zeros.
The two-change bound shows that every strict sublevel set of the lower branch
is an interval. Analyticity and the endpoint inequalities make its minimum
unique; the multiplicity version of that bound makes the minimum
nondegenerate. The endpoint velocity follows from
Corollary~\ref{lerchcont:regularity}. Full details, including the power-series
variation proof, are in \cite{lerchcont2026}.
\end{proof}

\begin{theorem}[All-branch eventual monotonicity]\label{lerchcont:monotonicity}
Fix $n\ge1$. For all sufficiently large $k$, every ordered zero
$a_{n,j,k}^{\rho}$ is strictly decreasing throughout $0\le\rho\le1$.
If $Q_n(x_{n,j})=0$, $c_{n,j}=e^{-x_{n,j}}$, and
$E_{n,j}=e^{1/c_{n,j}}$, then uniformly in $\rho$,
\[
 \partial_\rho a_{n,j,k}^{\rho}
 =-\frac{E_{n,j}}{(E_{n,j}-\rho)^2}+O_n(k^{-1}).
\]
\end{theorem}
The proof strengthens the existing gamma-concentration expansion to one
$\rho$ derivative, using
$\partial_\rho^qW_\rho(x)=q!e^{-qx}/(1-\rho e^{-x})^{q+1}$.
Its endpoint singularity has only fixed inverse-power growth, which the
large-shape gamma tails control. An explicit optimal threshold is not supplied.
Theorem~\ref{lerchcont:n2} shows why this eventual statement cannot be extended
to every derivative order.

\subsection{Small-parameter unfolding and an interior fold}
For fixed $n\ge2$ and all sufficiently small \emph{positive} $\rho$,
$F_{n,1}^{\rho}$ has one simple positive zero if $n$ is odd, and two if $n$
is even. The zero near $e^n$ persists analytically. When $n$ is even, the
other zero is
\[
 1-\left(\frac{(\log2)^{n-1}(n-\log2)}{4n}\right)^{1/(n-1)}
       \rho^{1/(n-1)}+O\bigl(\rho^{2/(n-1)}\bigr).
\]
This follows from expanding the elementary zero of multiplicity $n-1$ at
$a=1$. The endpoint $\rho=0$ itself must retain that multiplicity.

For $n=3$, there is a unique interior threshold $\rho_c$ at which a double
zero forms in $0<a<3/2$. There are no zeros in that interval for
$0<\rho<\rho_c$ and two simple zeros for $\rho_c<\rho\le1$. At the fold,
$F_\rho<0$ and $F_{aa}>0$, and the pair splits as
\[
 a_c\mp\sqrt{-2F_\rho/F_{aa}}\,\sqrt{\rho-\rho_c}+O(\rho-\rho_c).
\]
At the separate endpoint $\rho=0$, $a=1$ is a double zero.
Above the interior threshold, the total count is exactly three. Excluding an
additional outer pair at every intermediate $0<\rho<\rho_c$ remains a
separate conjecture, not a consequence of the localized fold proof.

\subsection{The exact fifth-index count}
\begin{theorem}[Sharp classical saturation at index five]\label{lerchcont:n5}
Every positive zero of $\gamma_5^{(k)}$ is simple, and its positive-zero count is
\[
 \#\{a>0:\gamma_5^{(k)}(a)=0\}=
 \begin{cases}3,&k=1,2,\\5,&k\ge3.\end{cases}
\]
\end{theorem}
\begin{proof}[Certified critical-point descent]
The verifier in \cite{lerchcont2026} gives five disjoint rational isolating
intervals for $F_{5,3}^{1}$. The global bound proves completeness and simplicity.
On those entire intervals, the predecessor $F_{5,2}^{1}$ has respective
certified signs $+,+,-,+,-$. Since its derivative is $-F_{5,3}^{1}$ and its
endpoint limits are $+\infty$ and $0^-$, it has exactly three simple zeros.
On isolating intervals for these three zeros, the next predecessor
$F_{5,1}^{1}$ has signs $-,+,-$. The same monotonic-interval argument proves
exactly three simple zeros at order one. Persistence from the saturated third
order gives five zeros at every higher order.

The calculation uses 35 exact sign enclosures with outward integer rounding,
rational logarithm bounds, and a proved Euler--Maclaurin remainder. In
particular, the predecessor signs are certified on whole critical-point
brackets, not at approximate critical points.
\end{proof}

The continuation settles the first $n\ge5$ small-order classical case left
open in the source, while preserving the existing eventual-saturation theorem.
It leaves the mixed $S_4$ identity and numerical period independence untouched.
